\documentclass[11pt]{article}
\usepackage{icmlko}
\begin{document}

\title{보폭은 잎에 남는다 --- 네 층을 파고 나서야 lr 을 찾았다}
\author{월드모델 실험실}
\date{노트 477 · 478 · 2026년 8월 3일}
\maketitle

\begin{abstract}
노트 475 가 물음 하나를 열어 놓고 못 닫았다 --- \emph{lr $0.16$ 은 학습
적합도를 $+0.0351$ 사는데 자 여섯 어디에서도 안 잃고 안 번다. 산 것이
어디로 갔나.} 노트 476 이 예보 층에서 찾다 실패했다(효과가 씨앗 잡음보다
작다). 이번엔 \textbf{모형을 열었다.}

\textbf{나무 구조에는 없다.} 잎/나무 $31.00$ · 깊이중앙 $11.0$ · 쓰인 축
$61/87$ --- 챔피언과 \texttt{it} $400$ 과 lr $0.16$ 이 \textbf{소수점까지
같다}. 축 사용 분포도 스피어만 $0.999$ 다.

\textbf{잎의 값에 있다.}

\begin{center}
\begin{tabular}{lrrr}
\hline
챔피언 대비 & 잎 수 & 잎당 평균 $|$값$|$ & 총 이동량 \\
\hline
lr $0.06$ & $\mathbf{1.000}$ & $\mathbf{0.664}$ & $0.664$ \\
lr $0.16$ & $\mathbf{1.000}$ & $\mathbf{1.425}$ & $1.425$ \\
\texttt{it} $400$ & $\mathbf{1.818}$ & $0.844$ & $1.535$ \\
\hline
\end{tabular}
\end{center}

\textbf{lr 은 잎을 한 개도 안 늘리고 값만 키운다. \texttt{it} 은 잎을
$1.818$배($=400/220$) 늘리고 각 잎은 오히려 작아진다.} 총 이동량은
$1.425$배와 $1.535$배로 비슷하고, 두 후보가 산 학습 적합도도
$+0.0351$과 $+0.0434$ 로 비슷하다.
\end{abstract}

\begin{figure}[t]
\centering
\includegraphics[width=\textwidth]{figs/f1.pdf}
\caption{\textbf{가} --- lr 은 세로로만, \texttt{it} 은 가로로만 움직인다.
점선은 총 이동량이 같은 곡선. \textbf{나} --- 네 층을 파고서야 찾았다.}
\end{figure}

\section{그래서 왜 판이 lr 을 못 보나}

\textbf{판은 순위 지표다.} 모든 잎을 같은 비율로 키우면 \emph{순서가
거의 안 바뀐다} --- 부스팅의 합이 단조 변환에 가깝게 늘어날 뿐이다.
그래서 lr $0.16$ 은 학습 적합도를 사고도 자 여섯 어디에도 안 나온다.
\emph{자가 모자란 게 아니라 살 것이 순위가 아니었다.}

\texttt{it} 은 다르다. 잎을 $1.818$배로 \textbf{늘려서} 총 이동량을
맞추는데, 새 잎들은 남은 잔차를 쫓으므로 \textbf{새로운 구별}을 만든다.
그 구별이 무리 \emph{사이} 쪽에 몰려서 --- 노트 $470 \cdot 473$ --- 판만
통과하고 나머지 다섯에서 잃는다.

\begin{quote}
\textbf{정리} --- 부스팅에서 lr 과 \texttt{it} 은 총 이동량을 같이
늘리지만 \textbf{lr 은 있던 구별을 키우고 \texttt{it} 은 새 구별을
만든다.} 순위 지표는 앞엣것에 둔감하고 뒤엣것에 민감하다.
\end{quote}

노트 461 이 이 둘을 \emph{배우는 용량}과 \emph{외우는 용량}이라 불렀다.
이제 그 말의 기계가 보인다 --- 외우기는 \textbf{새 잎}으로 한다.

\section{등록은 내 띠가 틀려 반증됐다}

미리 적은 것은 \emph{``lr 은 잎값을 lr 비율만큼 키우고 \texttt{it} 은
안 키운다''}였고 \texttt{it} 의 허용 띠를 $[0.9, 1.1]$ 로 잡았다.
\texttt{it} $400$ 의 잎값 배수는 $0.844$ 라 띠 밖이다. \textbf{반증으로
적는다.}

띠가 틀린 이유는 분명하다 --- 나무를 더 얹으면 \emph{각 나무가 맡을
잔차가 줄어들므로} 잎값이 작아진다. 그걸 안 생각하고 ``안 바뀐다''로
잡았다. lr 쪽도 비례가 아니다($0.16/0.10 = 1.60$ 인데 잎값은 $1.425$배)
--- 같은 이유다.

\section*{조항 --- 못 찾으면 층을 내린다}

\begin{quote}
\textbf{조항} --- 어떤 변경이 자에 안 잡히는데 학습 적합도는 움직이면,
\textbf{자를 더 만들기 전에 층을 내린다}: 점수 $\to$ 예보 $\to$ 구조
$\to$ 잎값. 이번에 넷째 층에서 찾았다. \emph{``자가 모자라다''는 마지막에
할 말이다.}
\end{quote}

\section*{남은 것}

\begin{itemize}
\item \textbf{잎값 자를 규약에 넣을까} --- 값싸고(라벨 안 씀) 뜻이
분명하다. 다만 이번처럼 \emph{순위에 안 나오는 변경}을 짚는 용도라
채택 판정에는 안 쓰인다. 기록으로만.
\item \textbf{계절이 제일 많이 쪼개지는 축이다}($0.0907$ · 노트 477).
노트 447 이 넣은 축인데 두 해 묵은 축들보다 더 쓰인다 --- 왜 그런지
안 봤다.
\item \textbf{원핫 몫이 챔피언 $0.0073$ 인데 후보 둘 다 $0.0056$} 이다.
절대 수로는 \texttt{it} $400$ 이 $1.39$배 더 쓴다 --- 노트 462 의
원핫 의존도(예보 기반)가 \texttt{it} 과 함께 오른 것과 맞는다.
\end{itemize}

\end{document}
